\documentclass[aspectratio=169,12pt]{beamer}
% Build from the repository root: make lecture-05-slides
\usepackage[T1]{fontenc}
\usepackage{lmodern}
\usepackage{amsmath,amssymb,booktabs,array}
\usepackage{pgfplots}
\pgfplotsset{compat=1.18}
\usefonttheme{professionalfonts}
\definecolor{Ink}{HTML}{16324F}
\definecolor{Blue}{HTML}{1F4E79}
\definecolor{Orange}{HTML}{C46A13}
\definecolor{Teal}{HTML}{007D80}
\definecolor{Muted}{HTML}{526477}
\setbeamercolor{normal text}{fg=Ink,bg=white}
\setbeamercolor{frametitle}{fg=Ink}
\setbeamercolor{structure}{fg=Blue}
\setbeamercolor{footline}{fg=Muted}
\setbeamerfont{frametitle}{size=\large,series=\bfseries}
\setbeamerfont{footline}{size=\tiny}
\setbeamersize{text margin left=8mm,text margin right=8mm}
\setbeamertemplate{navigation symbols}{}
\setbeamertemplate{headline}{}
\setbeamertemplate{frametitle}{\vspace{3mm}\insertframetitle\par\vspace{1mm}}
\newcommand{\conspectsection}{5.1}
\setbeamertemplate{footline}{%
  \hspace{8mm}\usebeamerfont{footline}\usebeamercolor[fg]{footline}%
  Conspect \S\,\conspectsection\hfill\insertframenumber\hspace{8mm}\vspace{3mm}}
\setlength{\parskip}{0.5em}
\setlength{\abovedisplayskip}{0.55em}
\setlength{\belowdisplayskip}{0.55em}
\newcommand{\Prob}{\mathbb{P}}
\newcommand{\Expect}{\mathbb{E}}
\DeclareMathOperator{\Var}{Var}
\DeclareMathOperator*{\argmax}{arg\,max}
\DeclareMathOperator{\MSE}{MSE}
\newcommand{\takeaway}[1]{\vfill{\small\color{Muted}#1\par}}
\newcommand{\plotcaption}[1]{\par\smallskip{\small #1\par}}
\setbeamertemplate{blocks}[default]
\setbeamerfont{block title}{size=\normalsize,series=\bfseries}
\setbeamerfont{block body}{size=\normalsize}
\newenvironment{coursebox}[2]{%
  \begingroup
  \setbeamercolor{block title}{fg=#1,bg=#1!13!white}
  \setbeamercolor{block body}{fg=Ink,bg=#1!4!white}
  \begin{block}{#2}
  \setlength{\parskip}{0.2em}
  \setlength{\abovedisplayskip}{0.3em}
  \setlength{\belowdisplayskip}{0.3em}
}{\end{block}\endgroup}
\newenvironment{definitionbox}[1]
  {\begin{coursebox}{Blue}{Definition: #1}}{\end{coursebox}}
\newenvironment{resultbox}[1]
  {\begin{coursebox}{Teal}{Result: #1}}{\end{coursebox}}
\newcommand{\EstDataPath}{../../figures/data}
\newcommand{\EstPlotWidth}{12cm}
\newcommand{\EstPlotHeight}{3.6cm}
% Shared analytic/simulation data charts for the conspect and the Beamer deck.
% Callers may override dimensions and the data path before loading this file.
\providecommand{\EstDataPath}{figures/data}
\providecommand{\EstPlotWidth}{11.8cm}
\providecommand{\EstPlotHeight}{4.8cm}
\definecolor{EstBlue}{HTML}{1F4E79}
\definecolor{EstOrange}{HTML}{C46A13}
\definecolor{EstTeal}{HTML}{007D80}
\definecolor{EstMuted}{HTML}{526477}
\pgfplotsset{est axis/.style={
  scale only axis,width=\EstPlotWidth,height=\EstPlotHeight,
  axis lines=left,enlargelimits=false,
  grid=major,grid style={black!9},
  tick label style={font=\footnotesize},label style={font=\footnotesize},
  legend style={draw=none,fill=none,font=\footnotesize,
    at={(0.5,1.05)},anchor=south,legend columns=-1,column sep=0.8em},
  every axis plot/.append style={no marks,line width=1.15pt},
}}

\newcommand{\EstLikelihoodPlot}{%
\begin{tikzpicture}\begin{axis}[est axis,
  xmin=0,xmax=1.2,ymin=0,ymax=1.06,
  xlabel={Candidate rate $\lambda$ (per minute)},ylabel={Relative likelihood},
  xtick={0,0.2,0.4,0.6,0.8,1,1.2}]
  \addplot[EstBlue] table[x=rate,y=n5] {\EstDataPath/ch5-likelihood.dat};
  \addlegendentry{$n=5$, total wait $15$}
  \addplot[EstOrange,dashed] table[x=rate,y=n50] {\EstDataPath/ch5-likelihood.dat};
  \addlegendentry{$n=50$, total wait $150$}
  \addplot[EstMuted,densely dotted,forget plot] coordinates {(0.333333,0) (0.333333,1)};
\end{axis}\end{tikzpicture}%
}

\newcommand{\EstUniformPlot}{%
\begin{tikzpicture}\begin{axis}[est axis,
  xmin=0,xmax=2,ymin=0,ymax=5.4,xlabel={Estimated upper endpoint},ylabel={Density},
  xtick={0,0.5,1,1.5,2}]
  \addplot[EstBlue,const plot] table[x=endpoint,y=moments] {\EstDataPath/ch5-uniform-estimates.dat};
  \addlegendentry{Moment estimator $2\bar X$}
  \addplot[EstOrange,const plot] table[x=endpoint,y=mle] {\EstDataPath/ch5-uniform-estimates.dat};
  \addlegendentry{MLE $\max_iX_i$}
  \addplot[EstMuted,densely dotted,forget plot] coordinates {(1,0) (1,5.2)};
  \node[anchor=north west,font=\footnotesize,text=EstMuted,fill=white,inner sep=2pt]
    at (axis cs:1.06,4.95) {True endpoint: $1$};
\end{axis}\end{tikzpicture}%
}

\newcommand{\EstVarianceLossPlot}{%
\begin{tikzpicture}\begin{axis}[est axis,
  xmin=0.5,xmax=6,ymin=-1,ymax=8.2,
  xlabel={Predicted standard deviation $\sigma$},ylabel={Loss contribution},
  xtick={1,2,3,4,5,6}]
  \addplot[EstBlue] table[x=sigma,y=total] {\EstDataPath/ch5-variance-loss.dat};
  \addlegendentry{Total: $\log\sigma+2/\sigma^2$}
  \addplot[EstOrange,dashed] table[x=sigma,y=scaled_residual] {\EstDataPath/ch5-variance-loss.dat};
  \addlegendentry{$2/\sigma^2$}
  \addplot[EstTeal,densely dotted] table[x=sigma,y=log_scale] {\EstDataPath/ch5-variance-loss.dat};
  \addlegendentry{$\log\sigma$}
  \addplot[EstBlue,only marks,mark=*,mark size=2pt,forget plot] coordinates {(2,1.19314718)};
  \node[font=\footnotesize,anchor=north,text=EstBlue,fill=white,inner sep=2pt]
    at (axis cs:2,0.8) {Minimum at $\sigma=2$};
\end{axis}\end{tikzpicture}%
}

\newcommand{\EstBetaPlot}{%
\begin{tikzpicture}\begin{axis}[est axis,
  xmin=0,xmax=1,ymin=0,ymax=3.4,
  xlabel={Bernoulli probability $p$},ylabel={Density over $p$},
  xtick={0,0.2,0.4,0.6,0.8,1}]
  \addplot[EstMuted,dashed] table[x=p,y=prior] {\EstDataPath/ch5-beta.dat};
  \addlegendentry{Prior: Beta$(2,2)$}
  \addplot[EstBlue] table[x=p,y=posterior] {\EstDataPath/ch5-beta.dat};
  \addlegendentry{Posterior: Beta$(9,5)$}
  \addplot[EstOrange,densely dotted,forget plot] coordinates {(0.6666667,0) (0.6666667,3.15)};
\end{axis}\end{tikzpicture}%
}

\newcommand{\EstGaussianPriorPlot}{%
\begin{tikzpicture}\begin{axis}[est axis,
  xmin=-3,xmax=5,ymin=0,ymax=0.61,
  xlabel={Candidate mean $\mu$},ylabel={Density / rescaled likelihood},
  xtick={-2,0,0.8,1.6,3,5},xticklabels={$-2$,$0$,$0.8$,$1.6$,$3$,$5$}]
  \addplot[EstMuted,dashed] table[x=mu,y=prior] {\EstDataPath/ch5-gaussian-prior.dat};
  \addlegendentry{Prior}
  \addplot[EstOrange,dashdotted] table[x=mu,y=likelihood] {\EstDataPath/ch5-gaussian-prior.dat};
  \addlegendentry{Rescaled likelihood}
  \addplot[EstBlue] table[x=mu,y=posterior] {\EstDataPath/ch5-gaussian-prior.dat};
  \addlegendentry{Posterior}
\end{axis}\end{tikzpicture}%
}

\title{Statistical estimation}
\subtitle{Lecture 5}
\author{}
\date{Scientific Computing 2026}
\hypersetup{pdftitle={Lecture 5: Statistical estimation},
  pdfsubject={Scientific Computing 2026, Chapter 5},pdfauthor={}}

\begin{document}

% 1
{\setbeamercolor{background canvas}{bg=Ink}
\begin{frame}[plain]
\color{white}
{\large Scientific Computing 2026\par}
\vfill
{\LARGE\bfseries Statistical estimation\par}
\vspace{0.7em}
{\large Method of moments, maximum likelihood, and MAP\par}
\vfill
{\large Lecture 5\par}\vspace{4mm}
\end{frame}}

% 2
\begin{frame}{Five observed waiting times}
Independent waits between requests, in minutes:
\[
  1,\ 2,\ 3,\ 4,\ 5.
  \qquad \bar t=3,\qquad \sum_i t_i=15.
\]
Keep the exponential model, but let its rate be unknown:
\[
  T_i\sim\operatorname{Exp}(\lambda),\qquad\lambda>0.
\]
Which value of $\lambda$ should we report?
\takeaway{Chapter 4 studied variation under a specified model.
Here the observations help us choose its parameter.}
\end{frame}

% 3
\begin{frame}{A statistical model}
\begin{definitionbox}{Statistical model}
A family of possible observation laws
\[
  \{P_\theta:\theta\in\Theta\}.
\]
The parameter space $\Theta$ specifies the allowed values.
\end{definitionbox}
Exponential rate: $\lambda\in(0,\infty)$.\\
Gaussian mean and variance: $(\mu,v)\in\mathbb R\times(0,\infty)$.
\takeaway{The family is an assumption about the observations.
Estimating its parameters does not verify that assumption.}
\end{frame}

% 4
\begin{frame}{Estimator, estimate, and identifiability}
\begin{definitionbox}{Estimator and estimate}
An estimator is a data-based rule $\widehat\theta=T(X_1,\ldots,X_n)$.
Its realized value $T(x_1,\ldots,x_n)$ is an estimate.
\end{definitionbox}
For the waits: the rule $1/\bar T_n$ gives $1/3$ per minute.
\vfill
\textbf{Identifiability:} different parameter values must give different laws.
For $X\sim\mathcal N(a^2,1)$, the data cannot distinguish $a$ from $-a$.
\end{frame}

\renewcommand{\conspectsection}{5.2}
% 5
\begin{frame}{Method of moments}
\begin{definitionbox}{Method of moments}
Match selected model moments to their empirical counterparts:
\[
  \Expect_{\widehat\theta}[X^j]=\frac1n\sum_{i=1}^n x_i^j.
\]
Solve for the unknown parameters within their allowed space.
\end{definitionbox}
For exponential waits,
\[
  \Expect_\lambda[T]=\frac1\lambda
  \quad\Longrightarrow\quad
  \widehat\lambda_{\mathrm{MM}}=\frac1{\bar t}=\frac13.
\]
\end{frame}

% 6
\begin{frame}{Two moments for a Gamma model}
Shape $\alpha>0$, rate $\beta>0$:
\[
  \Expect[X]=\frac\alpha\beta,\qquad
  \Var(X)=\frac\alpha{\beta^2}.
\]
With $v_n=n^{-1}\sum_i(x_i-\bar x)^2$, match
\[
  \bar x=\frac{\widehat\alpha}{\widehat\beta},\qquad
  v_n=\frac{\widehat\alpha}{\widehat\beta^2}.
\]
\begin{resultbox}{Gamma moment estimates}
\[
  \widehat\alpha=\frac{\bar x^2}{v_n},\qquad
  \widehat\beta=\frac{\bar x}{v_n},\qquad v_n>0.
\]
\end{resultbox}
\end{frame}

% 7
\begin{frame}{Matching the mean and spread}
For $1,2,3,4,5$ minutes,
\[
  \bar x=3,\qquad v_n=2,\qquad
  \widehat\alpha=4.5,\qquad\widehat\beta=1.5\text{ per minute}.
\]
The Gamma fit matches both summaries.\\
The exponential fit has variance $\bar x^2=9$.
\vfill
The denominator $n$ follows from matching raw moments:
\[
  v_n=\widehat m_2-\widehat m_1^2.
\]
\takeaway{If $v_n=0$, these Gamma moment equations have no finite positive solution.}
\end{frame}

% 8
\begin{frame}{Conditions for moment estimation}
The selected population moments must exist and distinguish the parameters.
\vfill
The LLN makes empirical moments approach population moments.
A well-defined, continuous inverse then makes the parameter estimates converge.
\vfill
Finite samples can still give no admissible solution, multiple solutions,
or estimates that react sharply to small changes in the moments.
\end{frame}

\renewcommand{\conspectsection}{5.3}
% 9
\begin{frame}{Likelihood}
\begin{definitionbox}{Likelihood for an i.i.d. sample}
With the observations held fixed,
\[
  L(\theta;x)=\prod_{i=1}^n f_\theta(x_i).
\]
Vary $\theta$ to compare candidate models for the recorded sample.
\end{definitionbox}
Discrete data use probability masses. Continuous data use densities.
\takeaway{Likelihood is not a distribution over $\theta$.
Dependent observations require their actual joint law.}
\end{frame}

% 10
\begin{frame}{Maximum likelihood and logarithms}
\begin{definitionbox}{Maximum likelihood estimator}
\[
  \widehat\theta_{\mathrm{ML}}\in\argmax_{\theta\in\Theta}L(\theta;x).
\]
The maximizer may be nonunique or may fail to exist.
\end{definitionbox}
The increasing logarithm preserves the maximum:
\[
  \ell(\theta;x)=\sum_i\log f_\theta(x_i).
\]
\takeaway{Drop only terms independent of the parameter.
Adding logs also avoids multiplying many very small numbers.}
\end{frame}

% 11
\begin{frame}{The exponential rate MLE}
For total observed waiting time $s>0$,
\[
  L(\lambda)=\lambda^n e^{-\lambda s},\qquad
  \ell(\lambda)=n\log\lambda-\lambda s.
\]
\[
  \ell'(\lambda)=\frac n\lambda-s,\qquad
  \ell''(\lambda)=-\frac n{\lambda^2}<0.
\]
\begin{resultbox}{Unique maximum}
\[
  \widehat\lambda_{\mathrm{ML}}=\frac n{s}=\frac1{\bar t}.
\]
\end{resultbox}
\end{frame}

% 12
\begin{frame}{Likelihood distinguishes candidate rates}
\centering\EstLikelihoodPlot
\plotcaption{Two illustrative datasets have average wait three minutes.
Both estimates are $1/3$ per minute. The vertical scale is $L(\lambda)/L(\widehat\lambda)$.}
\end{frame}

% 13
\begin{frame}{A parameter-dependent support}
For $X_i\sim\operatorname{Uniform}(0,\theta)$, let $x_{(n)}=\max_i x_i>0$.
\[
  L(\theta)=\theta^{-n}\mathbf1_{\{\theta\ge x_{(n)}\}}.
\]
An endpoint below an observation has likelihood zero.\\
Above the sample maximum, likelihood decreases.
\begin{resultbox}{Uniform endpoint MLE}
\[
  \widehat\theta_{\mathrm{ML}}=x_{(n)}.
\]
\end{resultbox}
\takeaway{The maximum lies at a boundary. No derivative-zero equation finds it.}
\end{frame}

% 14
\begin{frame}{Moment estimates and MLEs can differ}
Uniform model, with observations $0.2,0.4,0.5,0.7,0.9$:
\[
  \widehat\theta_{\mathrm{MM}}=2\bar x=1.08,\qquad
  \widehat\theta_{\mathrm{ML}}=\max_i x_i=0.9.
\]
\vfill
Moment matching uses $\Expect[X]=\theta/2$.\\
Likelihood also uses the support restriction.
\takeaway{Parameter constraints and global behavior matter.
A stationary point alone does not establish a maximum.}
\end{frame}

\renewcommand{\conspectsection}{5.4}
% 15
\begin{frame}{A Gaussian mean estimate}
For i.i.d. $X_i\sim\mathcal N(\mu,v)$ with $v>0$,
\[
  \ell(\mu,v)=-\frac n2\log(2\pi)-\frac n2\log v
  -\frac1{2v}\sum_i(x_i-\mu)^2.
\]
For fixed $v$, the identity
\[
  \sum_i(x_i-\mu)^2
  =\sum_i(x_i-\bar x)^2+n(\mu-\bar x)^2
\]
puts the maximum at $\widehat\mu=\bar x$.
\takeaway{Here $v$ denotes variance, not standard deviation.}
\end{frame}

% 16
\begin{frame}{A Gaussian variance estimate}
Write $Q=\sum_i(x_i-\bar x)^2>0$.
\[
  \frac{\partial\ell(\bar x,v)}{\partial v}
  =\frac{Q-nv}{2v^2}.
\]
The derivative changes from positive to negative at $v=Q/n$.
\begin{resultbox}{Gaussian MLEs}
\[
  \widehat\mu=\bar x,\qquad
  \widehat v=\frac1n\sum_i(x_i-\bar x)^2.
\]
\end{resultbox}
\takeaway{For identical observations, $Q=0$ and the likelihood grows without
bound as $v$ approaches zero.}
\end{frame}

\renewcommand{\conspectsection}{5.5}
% 17
\begin{frame}{Bias and mean squared error}
Repeat the entire sampling experiment with true parameter $\theta$ fixed.
\begin{definitionbox}{Bias and MSE}
\[
  \operatorname{Bias}_\theta(\widehat\theta)
  =\Expect_\theta[\widehat\theta]-\theta,
\]
\[
  \MSE_\theta(\widehat\theta)
  =\Expect_\theta[(\widehat\theta-\theta)^2].
\]
Unbiasedness means zero bias for every allowed parameter value.
\end{definitionbox}
\takeaway{Bias describes repeated-sample averages. One realized estimate has an error.}
\end{frame}

% 18
\begin{frame}{Bias and variance contribute to error}
Write
\[
  \widehat\theta-\theta
  =\bigl(\widehat\theta-\Expect_\theta\widehat\theta\bigr)
   +\bigl(\Expect_\theta\widehat\theta-\theta\bigr).
\]
Square and take expectations. The cross term has expectation zero.
\begin{resultbox}{Bias--variance decomposition}
\[
  \MSE_\theta(\widehat\theta)
  =\Var_\theta(\widehat\theta)+\operatorname{Bias}_\theta(\widehat\theta)^2.
\]
\end{resultbox}
\takeaway{Under squared error, smaller variance can compensate for some bias.}
\end{frame}

% 19
\begin{frame}{The $n-1$ variance correction}
For i.i.d. observations with finite variance,
\[
  \sum_i(X_i-\bar X)^2
  =\sum_i(X_i-\mu)^2-n(\bar X-\mu)^2.
\]
Taking expectations gives $n\sigma^2-n(\sigma^2/n)=(n-1)\sigma^2$.
\begin{resultbox}{Unbiased sample variance, $n\ge2$}
\[
  \Expect[\widehat v_{\mathrm{ML}}]=\frac{n-1}{n}\sigma^2,
  \qquad\Expect\!\left[\frac1{n-1}\sum_i(X_i-\bar X)^2\right]=\sigma^2.
\]
\end{resultbox}
\takeaway{The expectation identity needs no Gaussian assumption.
The Gaussian likelihood still chooses denominator $n$.}
\end{frame}

% 20
\begin{frame}{Squared error for a uniform endpoint}
For five i.i.d. Uniform$(0,1)$ observations:
\par\smallskip
{\small\renewcommand{\arraystretch}{1.5}
\begin{tabular*}{\textwidth}{@{\extracolsep{\fill}}lll@{}}
\toprule
Estimator & Bias & Mean squared error\\
\midrule
$2\bar X$ & $0$ & $1/15\approx0.0667$\\
$\max_i X_i$ & $-1/6$ & $1/21\approx0.0476$\\
\bottomrule
\end{tabular*}}
\vfill
The maximum's smaller spread outweighs its bias in this comparison.
\takeaway{For $M=\max_iX_i$, use $\Prob(M\le m)=m^5$ on $[0,1]$.
Integrating its density gives the moments used in the table.}
\end{frame}

% 21
\begin{frame}{The distribution of the estimates}
\centering\EstUniformPlot
\plotcaption{50,000 independent datasets of size five from Uniform$(0,1)$.
The MLE always underestimates the endpoint, but has smaller MSE here.}
\end{frame}

% 22
\begin{frame}{Consistency}
\begin{definitionbox}{Consistent estimator sequence}
For every $\varepsilon>0$,
\[
  \Prob_\theta(|\widehat\theta_n-\theta|>\varepsilon)\longrightarrow0.
\]
\end{definitionbox}
For exponential waits, the LLN gives $\bar T_n\to1/\lambda$.
Continuity of the reciprocal at this positive limit gives
$1/\bar T_n\to\lambda$ in probability.
\takeaway{This estimator is consistent but biased:
$\Expect[\widehat\lambda]=n\lambda/(n-1)$ for $n>1$.
Using only $X_1$ can be unbiased for a mean but inconsistent.}
\end{frame}

\renewcommand{\conspectsection}{5.6}
% 23
\begin{frame}{Gaussian regression gives squared error}
Fix the features and assume conditionally independent responses:
\[
  Y_i\mid x_i\sim\mathcal N(f_\theta(x_i),\sigma^2),
  \qquad\sigma^2>0\text{ known and shared}.
\]
\begin{resultbox}{Negative conditional log-likelihood}
\[
  -\ell(\theta)=\frac n2\log(2\pi\sigma^2)
  +\frac1{2\sigma^2}\sum_i(y_i-f_\theta(x_i))^2.
\]
\end{resultbox}
The parameter-dependent term is a positive multiple of squared error.
\end{frame}

% 24
\begin{frame}{Fitting a shared noise variance}
When the common variance $v$ is unknown, retain its logarithm:
\[
  -\ell(\theta,v)=\frac n2\log v
    +\frac{\operatorname{RSS}(\theta)}{2v}+\text{constant}.
\]
For fixed $\theta$ with positive residual sum of squares,
\[
  \widehat v(\theta)=\frac{\operatorname{RSS}(\theta)}{n}.
\]
\takeaway{A perfect fit has RSS $=0$ and allows variance collapse.
A finite fit then needs an additional constraint or a different model.}
\end{frame}

% 25
\begin{frame}{Predicting a different variance for each input}
Let $\mu_i=\mu_\theta(x_i)$ and $v_i=v_\theta(x_i)>0$.
\begin{resultbox}{Gaussian loss with input-dependent variance}
\[
  -\ell(\theta)=\frac12\sum_i\left[\log v_i+
    \frac{(y_i-\mu_i)^2}{v_i}\right]+\text{constant}.
\]
\end{resultbox}
A larger variance reduces the weighted residual but increases $\log v_i$.
\takeaway{$v_i=e^{s_i}$ enforces positivity.
It does not prevent collapse at a perfectly fitted observation.}
\end{frame}

% 26
\begin{frame}{The variance term has a cost}
\centering\EstVarianceLossPlot
\plotcaption{A fixed residual $y-\mu=2$ gives loss $\log\sigma+2/\sigma^2$.
Omitting $\log\sigma$ would reward arbitrarily large predicted noise.}
\end{frame}

% 27
\begin{frame}{Bernoulli responses give binary log loss}
For $Y_i\in\{0,1\}$, predict $p_i=p_\theta(x_i)\in(0,1)$.
\begin{resultbox}{Negative Bernoulli log-likelihood}
\[
  -\ell(\theta)=-\sum_i\bigl[y_i\log p_i+(1-y_i)\log(1-p_i)\bigr].
\]
\end{resultbox}
For a positive response:
\[
  -\log0.9\approx0.105,\qquad-\log0.1\approx2.303.
\]
Confident wrong predictions incur a large loss.
\end{frame}

% 28
\begin{frame}{The model, the objective, and its optimization}
The observation model determines the likelihood and training objective.
\vfill
An optimization algorithm searches for parameters with a small objective.
A stationary point need not be a global optimum.
\vfill
Training likelihood describes the fitted sample. Predictive evaluation
requires independent data and the assumptions discussed in Chapter 4.
\end{frame}

\renewcommand{\conspectsection}{5.7}
% 29
\begin{frame}{A prior distribution over the parameter}
Specify prior density $\pi(\theta)$ before using the current observations.
The model supplies $L(\theta;x)$ conditional on each candidate parameter.
\[
  \pi(\theta\mid x)=
  \frac{L(\theta;x)\pi(\theta)}{\int_\Theta L(u;x)\pi(u)\,du}.
\]
The denominator normalizes the posterior and is constant in $\theta$.
\takeaway{The denominator must be positive and finite.
For a discrete parameter, replace the integral by a sum.}
\end{frame}

% 30
\begin{frame}{Maximum a posteriori estimation}
\begin{definitionbox}{MAP estimate}
Maximize the posterior density, or posterior mass for a discrete parameter:
\[
  \widehat\theta_{\mathrm{MAP}}\in
  \argmax_{\theta\in\Theta}\{\ell(\theta;x)+\log\pi(\theta)\}.
\]
\end{definitionbox}
Equivalently, minimize negative log-likelihood plus negative log-prior.
\takeaway{MAP summarizes the posterior by its mode.
It does not report the whole posterior or its uncertainty.}
\end{frame}

% 31
\begin{frame}{A prior for a Bernoulli probability}
The Beta family places a density on $0<p<1$:
\[
  \pi(p)=\frac{p^{a-1}(1-p)^{b-1}}{B(a,b)},\qquad a,b>0,
\]
\[
  B(a,b)=\int_0^1u^{a-1}(1-u)^{b-1}\,du.
\]
The chosen $a,b$ describe the prior.\\
Beta$(2,2)$ has density $6p(1-p)$ and favors values near $1/2$.
\end{frame}

% 32
\begin{frame}{Updating a Beta prior}
For $k$ successes in $n$ independent Bernoulli trials,
\[
  L(p)\pi(p)\ \propto\ p^{k+a-1}(1-p)^{n-k+b-1}.
\]
Thus $p\mid x\sim\operatorname{Beta}(a+k,b+n-k)$.
\begin{resultbox}{Interior posterior mode}
\[
  \widehat p_{\mathrm{MAP}}=\frac{k+a-1}{n+a+b-2}.
\]
This formula requires $a+k>1$ and $b+n-k>1$.
\end{resultbox}
\end{frame}

% 33
\begin{frame}{Seven successes in ten trials}
\centering\EstBetaPlot
\plotcaption{Beta$(2,2)$ prior, Beta$(9,5)$ posterior.
MLE $=0.7$ and MAP $=2/3$; the dotted line marks the posterior mode.}
\end{frame}

% 34
\begin{frame}{Posterior mode and posterior mean}
For the Beta$(9,5)$ posterior,
\[
  \text{MAP}=\frac23\approx0.667,\qquad
  \Expect[p\mid x]=\frac9{14}\approx0.643.
\]
The mode maximizes density. The mean minimizes posterior expected
squared error, by the conditional-mean result from Chapter 3.
\takeaway{For continuous parameters, MAP depends on the parameterization:
transforming a density introduces a Jacobian factor.}
\end{frame}

% 35
\begin{frame}{A Gaussian prior for an unknown mean}
Observation model: $X_i\sim\mathcal N(\mu,\sigma^2)$, known $\sigma^2>0$.
Prior: $\mu\sim\mathcal N(m_0,\tau^2)$, with $\tau^2>0$.
\[
  -\log\pi(\mu\mid x)=
  \frac1{2\sigma^2}\sum_i(x_i-\mu)^2
  +\frac{(\mu-m_0)^2}{2\tau^2}+\text{constant}.
\]
\begin{resultbox}{Precision-weighted estimate}
\[
  \widehat\mu_{\mathrm{MAP}}=
  \frac{(n/\sigma^2)\bar x+(1/\tau^2)m_0}{n/\sigma^2+1/\tau^2}.
\]
\end{resultbox}
\end{frame}

% 36
\begin{frame}{How much weight the prior receives}
Precision means inverse variance. The sample mean has variance
$\sigma^2/n$, so its precision is $n/\sigma^2$.
\[
  n=4,\quad\sigma^2=4,\quad\bar x=1.6,\quad m_0=0,\quad\tau^2=1.
\]
Both precisions are one:
\[
  \widehat\mu_{\mathrm{ML}}=1.6,\qquad
  \widehat\mu_{\mathrm{MAP}}=0.8.
\]
The posterior is $\mathcal N(0.8,0.5)$.
\takeaway{Here the posterior mean and mode coincide.
With the prior fixed, more observations increase the relative data weight.}
\end{frame}

% 37
\begin{frame}{The Gaussian posterior}
\centering\EstGaussianPriorPlot
\plotcaption{The likelihood curve has been rescaled to integrate to one for comparison.
That rescaling does not make it a prior or posterior.}
\end{frame}

% 38
\begin{frame}{A Gaussian prior gives a quadratic penalty}
Gaussian responses with known variance $\sigma^2$ and prior
$w\sim\mathcal N(0,\tau^2I)$ give
\[
  -\log\pi(w\mid x,y)=\frac{\operatorname{RSS}(w)}{2\sigma^2}
  +\frac{\|w\|^2}{2\tau^2}+\text{constant}.
\]
\begin{resultbox}{Equivalent MAP objective}
\[
  \operatorname{RSS}(w)+\alpha\|w\|^2,\qquad
  \alpha=\frac{\sigma^2}{\tau^2}.
\]
\end{resultbox}
\takeaway{For a linear model, this is ridge regression.
The linear-algebra solution comes later in the course.}
\end{frame}

% 39
\begin{frame}{Loss normalization changes the penalty coefficient}
The same prior gives equivalent objectives
\[
  \sum_i(y_i-f_w(x_i))^2+\alpha\|w\|^2
\]
and
\[
  \frac1n\sum_i(y_i-f_w(x_i))^2+\frac\alpha n\|w\|^2.
\]
Replacing a sum by an average without adjusting the coefficient
changes the relative penalty strength.
\takeaway{An unpenalized intercept requires a different prior for that coefficient.}
\end{frame}

% 40
\begin{frame}{Three point-estimation rules}
\small\renewcommand{\arraystretch}{1.6}
\begin{tabular*}{\textwidth}{@{\extracolsep{\fill}}ll@{}}
\toprule
Method & Construction\\
\midrule
Method of moments & Match selected model and empirical moments\\
Maximum likelihood & Maximize the observation likelihood\\
MAP & Maximize likelihood times a specified prior\\
\bottomrule
\end{tabular*}\par\normalsize
\vfill
Their estimates can agree, but their definitions differ.
\takeaway{Point estimates alone do not quantify uncertainty.
Chapter 6 develops sampling distributions, confidence intervals, and bootstrap.}
\end{frame}

\end{document}
